% Generated draft: requires mathematical and source review.
\documentclass[11pt]{article}
\usepackage[margin=25mm]{geometry}
\usepackage{fontspec}
\setmainfont{Latin Modern Roman}
\usepackage{amsmath,amssymb,mathtools,unicode-math}
\setmathfont{Latin Modern Math}
\usepackage{enumitem,needspace,xurl,hyperref}
\setlength{\parindent}{0pt}
\setlength{\parskip}{5pt}
\setlength{\emergencystretch}{3em}
\newcommand{\N}{\mathbb{N}}
\newcommand{\Z}{\mathbb{Z}}
\newcommand{\Q}{\mathbb{Q}}
\newcommand{\R}{\mathbb{R}}
\newcommand{\C}{\mathbb{C}}
\newcommand{\F}{\mathbb{F}}
\newcommand{\A}{\mathbb{A}}
\newcommand{\PP}{\mathbb{P}}
\newcommand{\E}{\mathbb{E}}
\newcommand{\eps}{\varepsilon}
\newcommand{\dd}{\,\mathrm{d}}
\begin{document}
\section*{UnsolvedMath ID: 3100077}
\textbf{Source problem code:} catalogue entry
\subsection{Definitions and mathematical statement}
The supplied question asks whether, for every integer \(k>1\), there is a word unavoidable over a \(k\)-letter alphabet but not over a \((k-1)\)-letter alphabet. The meaning of ``word'' and ``unavoidable'' is not specified in the source, and the direction of comparison requires clarification.

For example, in the standard pattern interpretation, let \(P\) be a nonempty finite word over a finite alphabet of variables. For a finite target alphabet \(\Sigma\), an occurrence of \(P\) in a word \(w\) is a contiguous factor \(h(P)\), where \(h\) is a morphism assigning a nonempty word over \(\Sigma\) to each variable. Write \(U_k(P)\) if there is an integer \(L\ge1\) such that every word in \(\Sigma^L\), with \(|\Sigma|=k\), contains such an occurrence. Under this interpretation the source's literal quantified target is
\[
 \forall k\ge2\ \exists P:\quad U_k(P)\ \text{and}\ \neg U_{k-1}(P).
\]
However, the standard meaning of a word over \(\Sigma\) permits use of a subset of \(\Sigma\). Consequently \(U_k(P)\) implies \(U_{k-1}(P)\) by restricting the target alphabet. Thus the stated direction is incompatible with the ordinary monotonicity convention. Reversing the two alphabet sizes would change the source question and is not done here. A literal fixed-factor meaning of ``word'' also differs from the pattern interpretation and is not selected without evidence.

The record supplies neither the definition from its linked reference nor that link's destination. It therefore does not determine an unambiguous intended nontrivial problem. Its research summary reports no verified definitive answer and labels it open/unverified; the October 2020 source-page open label was not checked item by item against later literature. Those are source-status reports, not verification of a precise mathematical target.
\subsection{Short English statement}
The source asks for a word unavoidable over \(k\) letters but avoidable over \(k-1\) letters, for each \(k>1\). Under the usual pattern definition the implication goes the other way, so the source's meaning or ordering needs confirmation; no corrected question is silently substituted.
\subsection{Sources}
\begin{itemize}
\item Local registry: \url{data/unsolvedmath/problems.json}, permanent UnsolvedMath ID \(3100077\), source code catalogue entry; supplied statement, background, and research summary, updated 2026-08-06.
\item Cooper, \emph{Combinatorial Problems I Like}, October 2020 snapshot, Words and Codes, source-order bullet 77: \url{https://people.math.sc.edu/cooper/combprob.html}. HTML extraction accessed 2026-07-29.
\item The original statement says ``See this'', but the supplied record contains no destination or authoritative definition from that reference. The research summary mentions the unavoidable-pattern literature without an identifiable bibliographic citation.
\end{itemize}
\end{document}
